\chapter{اللقاء السادس: التفسير العملي لسورة الفاتحة}

\section{تأويل (الحمد لله)}
السلام عليكم ورحمة الله وبركاته. الحمد لله رب العالمين، والصلاة والسلام على نبينا محمد وعلى آله وصحبه أجمعين. وصلنا بحمد الله إلى بداية تفسير الإمام \rawi{الطبري} لسورة الفاتحة.

\isnad{قال أبو جعفر رحمه الله:} «القول في تأويل فاتحة الكتاب: \begin{ayah}الْحَمْدُ لِلَّهِ\end{ayah}. معنى (الحمد لله): الشكر لله خالصاً دون سائر ما يُعبد من دونه، ودون كل ما برأ من خلقه، بما أنعم على عباده من النعم التي لا يحصيها العدد...».

ثم ساق \rawi{الطبري} الأخبار، \isnad{فحدث عن} \rawi{ابن عباس} \isnad{أنه قال:} «الحمد هو الشكر لله والاستخذاء لله والاقرار بنعمته وهدايته وابتدائه وغير ذلك». \isnad{وحدث عن} \rawi{كعب الأحبار} \isnad{أنه قال:} «الحمد لله ثناء الله».

ونبه الشيخ هنا إلى أصل مهم في قراءة الكتاب: أننا في هذا المقام نطلب أولاً فهم اختيار \rawi{الطبري} نفسه، وكيف بنى قوله، وكيف استدل له، قبل الانتقال إلى الموازنة الواسعة بينه وبين غيره. فالمقصود في هذه المرحلة أن نفهم منطق الإمام من داخل كتابه.

\begin{fawaid}[تفسير الطبري ليس مجرد نقل للمأثور]
بعض الناس يصف تفسير الطبري بأنه "تفسير بالمأثور" ويقصد بذلك أنه مجرد ناقل للآثار؛ وهذا خطأ شنيع. الطبري يجتهد في الوصول إلى المعنى، ثم يذكر الآثار للاستدلال، ويورد الإشكالات اللغوية والعقلية ويجيب عنها. فتصنيفه للأقوال، وترجيحه بينها، ونقده لها، يدل على أنه إمام مجتهد محقق وليس مجرد جامع للمرويات.
\end{fawaid}

\subsection{دخول الألف واللام في (الحمد)}
طرح \rawi{الطبري} إشكالاً لغوياً: وما وجه إدخال الألف واللام في (الحمد)؟ وهلا قيل: حمداً لله رب العالمين؟
ثم أجاب بأن دخول الألف واللام يُنبئ عن أن معناه "جميع المحامد والشكر الكامل لله". ولو أسقطتا منه ما دل إلا على أن حمد قائل ذلك لله دون المحامد كلها. 
ولذلك تتابعت قراءة القراء وعلماء الأمة على (رفع) الدال من (الحمدُ) دلالة على الاستغراق والكمال.

\begin{fawaid}[الاحتجاج بالضعيف لبيان المعنى اللغوي]
قد يورد الطبري آثاراً فيها ضعف من جهة الإسناد (مثل بعض مرويات كعب الأحبار أو المراسيل والمقاطيع)، لا ليتعبد بتصحيحها حديثياً، بل لأن معناها اللغوي صحيح، فيستأنس بها في بيان دلالات الألفاظ، خاصة في الكلمات البينة كـ "الحمد".
\end{fawaid}

\begin{fawaid}[قد يقرر الطبري المعنى ثم يسوق له الآثار]
من الملامح الظاهرة في هذا الموضع أن \rawi{الطبري} قد يذكر المعنى الذي اختاره أولاً، ثم يتبعه بالآثار التي يرى أنها تشهد له. فلا يلزم أن يكون ترتيب الاستدلال عنده دائماً: أثر ثم نتيجة، بل قد يقرر الخلاصة أولاً ثم يسوق ما يعضدها.
\end{fawaid}

\separator

\section{تأويل (رب العالمين)}
\isnad{قال أبو جعفر:} «وأما تأويل قوله (رب)، فإن الرب في كلام العرب منصرف على معانٍ: السيد المطاع فيهم يُدعى رباً... والرجل المصلح يُدعى رباً... والمالك للشيء يُدعى ربه».

ثم استشهد \rawi{الطبري} بأشعار العرب كقول \rawi{لبيد بن ربيعة} و\rawi{النابغة الذبياني} و\rawi{الفرزدق} لإثبات هذه المعاني. وخلص إلى أن: «فربنا جل ثناؤه: السيد الذي لا شبه له في سؤدده، والمصلح أمر خلقه بما أسبغ عليهم من نعمه، والمالك الذي له الخلق والأمر».

\begin{fawaid}[تقديم الشواهد اللغوية على الآثار]
من منهج الطبري في تفسير الألفاظ أنه يبدأ غالباً بالشواهد الشعرية ودلالات الكلمة في لسان العرب، ثم من مجموع هذه الدلالات يخلص إلى المعنى، ثم يسوق بعد ذلك آثار الصحابة والتابعين المؤيدة له، ليجمع بين التفسير اللغوي والأثري؛ وذلك ليكون كتابه جامعاً مغنياً عما سبقه من كتب اللغويين (كالفراء وأبي عبيدة).
\end{fawaid}

\subsection{تأويل (العالمين)}
\isnad{قال الطبري:} «والعالمون جمع عالَم... والعالَم اسم لأصناف الأمم، وكل صنف منها عالَم».
ثم ساق آثاراً عن \rawi{ابن عباس} و\rawi{سعيد بن جبير} و\rawi{مجاهد} أن المراد بالعالمين: «الجن والإنس».

\separator

\section{تأويل (الرحمن الرحيم) ومسألة البسملة}
\isnad{قال أبو جعفر:} «قد مضى البيان عن تأويل قوله (الرحمن الرحيم) في تأويل (بسم الله الرحمن الرحيم)، فأغنى ذلك عن إعادته في هذا الموضع».

ثم أوضح \rawi{الطبري} رأيه في البسملة قائلاً: «إذ كنا لا نرى أن \textit{بسم الله الرحمن الرحيم} من فاتحة الكتاب آية». واحتج بأنها لو كانت آية، لكان تكرار قوله \begin{ayah}الرَّحْمَٰنِ الرَّحِيمِ\end{ayah} في الآية التي تليها تكراراً بلا فاصل، وهذا غير معهود في كتاب الله أن تتكرر آية بكمالها وبمعنى واحد متجاورة دون فاصل يفصل بينهما بكلام يخالف معناهما.

\separator

\section{تأويل (ملك يوم الدين) والترجيح بين القراءات}
\isnad{قال أبو جعفر:} «القراء مختلفون في تلاوة (مَالِكِ يَوْمِ الدِّينِ)؛ فبعضهم يتلوه \textit{مَلِكِ يوم الدين}، وبعضهم يتلوه \textit{مَالِكِ يوم الدين}، وبعضهم يتلوه \textit{مَالَكَ يوم الدين} بنصب الكاف».

أما القراءة بالنصب فقد اعتبرها شاذة مردودة. ووازن بين (مَلِك) و(مَالِك) ورجح الأولى قائلاً: «وأولى التأويلين بالآية وأصح القراءتين في التلاوة عندي: قراءة من قرأ (مَلِكِ) بمعنى المُلك؛ لأن في الإقرار له بالانفراد بالمُلك إيجاباً لانفراده بالمِلك، وفضيلة زيادة المُلك على المِلك... ولأن الله أخبر عن نفسه في الآية السابقة بأنه (رب العالمين) والرب يحوي معنى المالك».

\begin{fawaid}[الترجيح بين القراءات المتواترة]
الطبري قد يرجح قراءة صحيحة على قراءة أخرى صحيحة ويذكر علته في ذلك، وهو هنا رجح قراءة (مَلِك) لسببين:
الأول: أن (المَلِك) يتضمن معنى (المَالِك) وزيادة، فكل مَلِكٍ مالك وليس العكس.
الثاني: أن معنى (المالك) قد ثبت سلفاً بكلمة (رب)، فلو قرأنا (مالك) لكان تكراراً لمعنى تقدم، فقراءة (مَلِك) تفيد معنى جديداً للسامع.
\end{fawaid}

\subsection{معنى (يوم الدين)}
\isnad{قال الطبري:} «والدين في هذا الموضع بتأويل الحساب والمجازاة بالأعمال، كما قال \rawi{كعب بن جعيل}: \textit{إذا ما رمونا رميناهمُ *** ودِنّاهمُ مثل ما يقرضون}».

\separator

\section{تأويل (إياك نعبد وإياك نستعين) والرد على القدرية}
\isnad{قال أبو جعفر:} «وتأويل قوله (إياك نعبد): لك اللهم نخشع ونذل ونستكين... لأن العبودية عند جميع العرب أصلها الذلة، ومنه سمي الطريق المذلل الذي وطئته الأقدام معبداً».

وأما تأويل (وإياك نستعين) فهو طلب المعونة من الله على طاعته فيما بقي من العمر. وقد طرح \rawi{الطبري} إشكالاً للرد على القدرية المعتزلة:
\isnad{قال:} «فإن قال قائل: وما معنى أمر الله عباده بأن يسألوه المعونة على طاعته؟ أوجائز وقد أمرهم بطاعته ألا يعينهم عليها؟... قيل: وفي أمر الله عباده أن يقولوا (إياك نعبد وإياك نستعين) الدليل على فساد قول القائلين بالتفويض من أهل القدر الذين أحالوا أن يأمر الله أحداً بفرض إلا بعد إعطائه المعونة والقدرة...».

\begin{fawaid}[الرد على القدرية من (إياك نستعين)]
هذا الاستنباط العقدي يهدم أصول القدرية؛ فالله خلق في العبد قدرة ومشيئة ليحاسبه عليها، وإعانته للعبد على الطاعة "تفضلٌ منه" وليست "واجبة عليه". فلو كان العبد مستقلاً بفعله (كما يزعمون) لما كان هناك معنى لسؤال العبد ربه المعونة قائلاً (وإياك نستعين)، لأن المعونة على قولهم واجبة حتمية قبل الفعل!.
\end{fawaid}

\begin{fawaid}[لا يستغني العبد عن ربه في أصل الطاعة ولا في تمامها]
يبين هذا الموضع أن العبد وإن كان مخاطَباً مأموراً وله قدرة واختيار، فإنه لا يستقل بالطاعة عن ربه، بل يفتقر إلى هدايته وإعانته وتثبيته. فاجتماع التكليف مع الافتقار إلى المعونة من أظهر معاني العبودية.
\end{fawaid}

\separator

\section{تأويل (اهدنا الصراط المستقيم)}
\isnad{قال أبو جعفر:} «ومعنى قوله (اهدنا الصراط المستقيم) في هذا الموضع عندنا: وفقنا للثبات عليه... ومعناه نظير معنى قوله (وإياك نستعين) مسألة العبد ربه الثبات على العمل بطاعته وإصابة الحق».

\begin{fawaid}[تفسير الهداية بالتوفيق]
الهداية في القرآن لها دلالات: التعليم والإرشاد، والتوفيق والتثبيت. والطبري حصر معنى الهداية هنا في (التوفيق للثبات على الطاعة)، وقد ردّ شيخ الإسلام \rawi{ابن تيمية} هذا الحصر، وبيّن أن العبد محتاج إلى الهداية بمعانيها المتعددة مجتمعة (التعليم، وتحبيب الإيمان في القلب، والتوفيق للعمل، والتثبيت عليه).
\end{fawaid}

\isnad{قال الطبري في الصراط:} «أجمعت الحجة من أهل التأويل جميعاً على أن الصراط المستقيم هو الطريق الواضح الذي لا اعوجاج فيه، وكذلك في لغة جميع العرب... والذي هو أولى بتأويل الآية أن يكون معنياً به: وفقنا للثبات على ما ارتضيته... من الإسلام وتصديق الرسل والتمسك بالكتاب».

\separator

\section{تأويل (صراط الذين أنعمت عليهم غير المغضوب عليهم ولا الضالين)}
ساق \rawi{الطبري} آثار الصحابة القاضية بأن (الذين أنعمت عليهم) هم الملائكة والنبيون والصديقون والشهداء والصالحون. 

\subsection{توجيه الخفض في (غير)}
تكلم \rawi{الطبري} في إعراب (غيرِ) بخفض الراء، وذكر أنها إما نعت لـ (الذين)، أو خُفضت بنية تكرار العامل وهو (صراط)، كأنه قال: صراط الذين أنعمت عليهم، صراط غير المغضوب عليهم.

\begin{fawaid}[الاقتصار على الإعراب الخادم للمعنى]
نصّ الطبري على أنه لم يعترض في وجوه الإعراب في كتابه إلا لاضطرار الحاجة لتجلية المعنى وكشف وجه التأويل. وهذا يفصل بين منهجه كـ "مفسر" يطلب المعنى، وبين منهج اللغويين الذين يجعلون إعراب القرآن غاية في ذاته.
\end{fawaid}

\subsection{من هم المغضوب عليهم والضالون؟}
ساق \rawi{الطبري} الإجماع والأحاديث (كحديث \rawi{عدي بن حاتم}) على أن (المغضوب عليهم) هم اليهود، وأن (الضالين) هم النصارى.

وقد رد \rawi{الطبري} على القدرية الذين استدلوا بإضافة الضلال إلى العبد (الضالين) لنفي أن يكون الله هو من أضلهم، \isnad{فقال:} «وقد ظن بعض أهل الغباء من القدرية أن في وصف الله النصارى بالضلال... دلالة على صحة ما قاله إخوانه جهلاً منه بسعة كلام العرب... فكيف بالفعل الذي يكتسبه العبد كسباً، ويوجده الله عيناً ونشأة؟ فلذلك أحرى أن يضاف إلى مكتسبه كسباً، وإلى الله إيجاداً وإنشاءً».

\begin{fawaid}[خلق أفعال العباد والرد على القدرية]
من بديع نقض الطبري للقدرية توضيحه أن "إضافة الفعل للعبد" لا تنفي خلق الله له؛ فالضلال صفة اكتسبها العبد ويوصف بها، ولكن الله هو الذي خلقه وأوجده وأضله بقدره. فالعبد فاعل لفعله والله خالق لهذا الفعل.
\end{fawaid}

\separator

\section{صفة الغضب لله تعالى ونقد المتكلمين}
عند تفسير (المغضوب عليهم)، تطرق \rawi{الطبري} لمسألة عقدية جليلة، \isnad{فقال:} «واختُلِف في صفة الغضب من الله... فقال بعضهم: غضب الله... عقوبته (وهذا تأويل). وقال بعضهم: الغضب منه معنى مفهوم كالذي يُعرف من معاني الغضب، غير أنه وإن كان كذلك من جهة الإثبات فمخالف معناه منه معنى ما يكون من غضب الآدميين الذي يزعجهم ويحركهم... ولكنه له صفة، كما العِلم له صفة والقدرة له صفة على ما يُعقل من جهة الإثبات...».

\begin{fawaid}[إثبات صفات الأفعال لله تعالى]
هذا النص من أعظم نصوص الطبري في تقرير عقيدة السلف في الصفات (كالغضب والمجيء)؛ حيث أثبت أصل الصفة كما تُفهم في اللغة، ونفى عنها لوازم النقص البشري (كالانفعال والاضطراب)، وألزم المعطلة بالمتفق عليه، قائلاً: كما أنكم تثبتون لله العلم والقدرة دون أن يشبه علم وقدرة البشر، فكذلك أثبتوا الغضب والمحبة بلا تمثيل ولا تحريف.
\end{fawaid}

\separator

\section{الرد على الملاحدة في حكمة تطويل الفاتحة}
ختم \rawi{الطبري} سورة الفاتحة بدفع إيراد للملاحدة الذين يزعمون أن معاني الفاتحة يمكن اختصارها في (مالك يوم الدين، وإياك نعبد وإياك نستعين) فما الحاجة لتطويلها بالحمد والثناء؟

\isnad{فأجاب:} «إن الله جمع لنبينا صلى الله عليه وسلم ولأمته في القرآن معاني لم يجمعها في كتاب نزل قبله... فمنها نظمه العجيب، ورصفه الغريب... وبما فيه من تحميد وتمجيد وثناء عليه؛ تنبيه للعباد على عظمته وسلطانه ليذكروه بآلائه... فذلك وجه إطالة البيان في سورة أم القرآن».

وبهذا ينتهي التفسير العملي لسورة الفاتحة من كتاب الإمام \rawi{الطبري} رحمه الله.
